\documentclass[../DoAn.tex]{subfiles}
\begin{document}

Phụ lục này trình bày các biểu đồ use case phân rã bổ sung và đặc tả use case ``Thẩm định và lập báo cáo''. Các nội dung được chuyển khỏi Chương~\ref{chapter:Related_works} nhằm giữ chương chính tập trung vào chuỗi nghiệp vụ cốt lõi.

\section{Các biểu đồ use case phân rã bổ sung}

Các biểu đồ dưới đây chi tiết hóa những nhóm chức năng đã xuất hiện trong biểu đồ use case tổng quát. Hình~\ref{fig:appendix-account-use-case}, Hình~\ref{fig:appendix-product-use-case} và Hình~\ref{fig:appendix-wallet-use-case} mô tả tài khoản, sản phẩm và tài chính; Hình~\ref{fig:appendix-order-use-case}, Hình~\ref{fig:appendix-dispute-use-case} và Hình~\ref{fig:appendix-admin-use-case} mô tả đơn hàng, tranh chấp và quản trị. Tác nhân ``Thời gian'' được sử dụng để biểu diễn điều kiện kích hoạt tác vụ nền theo lịch, không phải một người dùng bên ngoài hệ thống.

\begin{figure}[htbp]
    \centering
    \includegraphics[width=\linewidth,height=0.72\textheight,keepaspectratio]{chapter-2/account-use-case.png}
    \caption{Biểu đồ use case phân rã nghiệp vụ quản lý tài khoản và hồ sơ}
    \label{fig:appendix-account-use-case}
    \caption*{\textit{Nguồn: Tổng hợp từ mã nguồn và phạm vi WoodCert Auction.}}
\end{figure}

\begin{figure}[p]
    \centering
    \includegraphics[width=\linewidth,height=0.72\textheight,keepaspectratio]{chapter-2/product-use-case.png}
    \caption{Biểu đồ use case phân rã nghiệp vụ quản lý sản phẩm và thẩm định}
    \label{fig:appendix-product-use-case}
    \caption*{\textit{Nguồn: Tổng hợp từ mã nguồn và phạm vi WoodCert Auction.}}
\end{figure}

\begin{figure}[p]
    \centering
    \includegraphics[width=\linewidth,height=0.72\textheight,keepaspectratio]{chapter-2/wallet-use-case.png}
    \caption{Biểu đồ use case phân rã nghiệp vụ quản lý ví và thanh toán}
    \label{fig:appendix-wallet-use-case}
    \caption*{\textit{Nguồn: Tổng hợp từ mã nguồn và phạm vi WoodCert Auction.}}
\end{figure}

\begin{figure}[p]
    \centering
    \includegraphics[width=\linewidth,height=0.72\textheight,keepaspectratio]{chapter-2/order-use-case.png}
    \caption{Biểu đồ use case phân rã nghiệp vụ xử lý đơn hàng và giao nhận}
    \label{fig:appendix-order-use-case}
    \caption*{\textit{Nguồn: Tổng hợp từ mã nguồn và phạm vi WoodCert Auction.}}
\end{figure}

\begin{figure}[p]
    \centering
    \includegraphics[width=\linewidth,height=0.72\textheight,keepaspectratio]{chapter-2/dispute-use-case.png}
    \caption{Biểu đồ use case phân rã nghiệp vụ giải quyết tranh chấp}
    \label{fig:appendix-dispute-use-case}
    \caption*{\textit{Nguồn: Tổng hợp từ mã nguồn và phạm vi WoodCert Auction.}}
\end{figure}

\begin{figure}[p]
    \centering
    \includegraphics[width=\linewidth,height=0.72\textheight,keepaspectratio]{chapter-2/admin-use-case.png}
    \caption{Biểu đồ use case phân rã nghiệp vụ quản trị hệ thống và tác vụ tự động}
    \label{fig:appendix-admin-use-case}
    \caption*{\textit{Nguồn: Tổng hợp từ mã nguồn và phạm vi WoodCert Auction.}}
\end{figure}

\FloatBarrier

\section{Đặc tả use case thẩm định và lập báo cáo}

Bảng~\ref{tab:appendix-uc-submit-appraisal-report} trình bày tác nhân, điều kiện, luồng xử lý và ngoại lệ của use case thẩm định và lập báo cáo.

\begin{table}[htbp]
    \centering
    \caption{Đặc tả use case: Thẩm định và lập báo cáo}
    \label{tab:appendix-uc-submit-appraisal-report}
    \footnotesize
    \setlength{\tabcolsep}{4pt}
    \renewcommand{\arraystretch}{1.05}
    \begin{tabular}{|>{\raggedright\arraybackslash}p{0.21\textwidth}|>{\raggedright\arraybackslash}p{0.69\textwidth}|}
        \hline
        \textbf{Mục} & \textbf{Nội dung} \\
        \hline
        Tác nhân chính & Thẩm định viên. \\
        \hline
        Tác nhân liên quan & Người bán. \\
        \hline
        Mục tiêu & Tiếp nhận yêu cầu, kiểm tra hiện vật và nộp báo cáo để xác định sản phẩm đạt hoặc không đạt điều kiện đấu giá. \\
        \hline
        Tiền điều kiện & Thẩm định viên có quyền xử lý; sản phẩm ở \texttt{PENDING\_APPRAISAL} hoặc \texttt{UNDER\_APPRAISAL}; quyền tiếp nhận còn hiệu lực; chưa có báo cáo cho sản phẩm. \\
        \hline
        Hậu điều kiện & Báo cáo được lưu một lần với mã chứng thư và mã băm SHA-256; ảnh chứng minh được gắn nếu có; sản phẩm chuyển sang \texttt{APPRAISED} hoặc \texttt{REJECTED}; quyền tiếp nhận được giải phóng. \\
        \hline
        Luồng chính & Thẩm định viên tiếp nhận yêu cầu; kiểm tra hiện vật theo quy trình ngoài phần mềm; nhập vật liệu xác minh, nguồn gốc, tình trạng, giá trị ước tính và kết quả; tùy chọn tải ảnh chứng minh; Hệ thống kiểm tra dữ liệu, tạo báo cáo và cập nhật trạng thái sản phẩm. \\
        \hline
        Ngoại lệ & Quyền tiếp nhận thuộc Thẩm định viên khác hoặc đã hết hạn; sản phẩm sai trạng thái; báo cáo đã tồn tại; kết quả không đạt nhưng thiếu lý do; ảnh không hợp lệ hoặc không thuộc người tải. \\
        \hline
        Ghi chú & Chức năng tra cứu trả về mã băm đã lưu nhưng chưa tính lại mã băm để xác minh tính toàn vẹn tại thời điểm tra cứu. Báo cáo không có luồng cập nhật hoặc xóa sau khi nộp trong ứng dụng hiện hành. \\
        \hline
    \end{tabular}
\end{table}

\end{document}
